% Online Resource 1, Section S3 (writer F, journal revision of 2026-10-13).
% Combines the design material moved from sections/supplement.tex of the
% rule-assembled attempt with the technical passages cut from its Methods.
% Moved sentences keep their wording except for terminology (influence for
% partner drift, contacts for partners, persistence model for null, signal
% share for response signal share, plain stage names before gate and tier
% codes) and American spelling. Claim markers Mxx and SUPxx keep their rows.

This section gives the design and inference details that Section~\ref{main-sec:methods} of the article summarizes. It joins material from the supplement of the manuscript assembled by the prespecified reading rules (Section~\ref{sec:rule-record}) with the model definitions and technical passages moved from the article's Methods. These editorial relocations do not change the protocol or the results. Codes of validation-gate stages (G0--G2c), inventory tiers and protocol versions are listed in the glossary (Section~\ref{sec:glossary}). The protocol history, cumulative exposure, validation gate, execution and the operational amendment of 1 October 2026 (v4.4) are recorded in Section~\ref{sec:record}; docking and the other implementation checks are in Section~\ref{sec:verification}.

\subsection{Model and channel mechanics}\label{sec:dd-channels}

% SUP01
The influence term is $D_{it}=8^{-1}\sum_{j\in\mathcal N_i}f(z_{jt}-z_{it})$, with $f(g)=g$ when $|g|\leq1$ and $f(g)=-0.25g$ otherwise, standardized across the population, to $d_{it}$, before entering the update. Zero variance of the influence term at positive strength raises an error and is not assigned a variance share. The frozen v4 code computes no diagnostic of the realized share of influence in latent change, so $a$ remains a nominal share; the exploratory variance-explained analysis (X2, Section~\ref{sec:exploratory}) reports that share outside the frozen protocol. Anchoring uses $R_N((1-\omega)h+\omega z_0)$ for $\omega\in\{0.25,0.50,0.75,1\}$; $\omega=1$ is a degenerate constant-latent case and is not declared identical to B. D uses $\rho(z)=\rho_0+\gamma(1-\rho_0)(1-e^{-z^2/2})$ inside the transported update. C multiplies response error by a normalized $1+2e^{-z^2/2}$ factor and heaps each response at zero with probability 0.10 before ordinal allocation. Response-style coefficients are 0.25 and 0.50; conditioning increments are 0.10 and 0.20; their strongest combination is exploratory. The same annual category target is used for competing arms. Matching shares masks within coupled pairs; test masks are independently drawn by class.

% M68 (moved from Methods)
Strength $a$ has no direct observable meaning. The frozen v1 protocol committed to report the realized partial $R^2$ of influence by strength, steps and $b$, but the frozen v4 code computes no such diagnostic. The v1 value measured at $b=0$ under the v1 law should not be transferred automatically to this design.

% M14 (moved from Methods)
Let $X_\phi$ start at $R_N(\epsilon_0)$ and apply $x_t=R_N(\phi x_{t-1}+\sqrt{1-\phi^2}\epsilon_t)$ at every step, including all 24 steps in that sensitivity. Persistence model B adds an independent, untransported iid standard-normal trait $T_i$. Parameters $\lambda$ and $\pi$ are pre-transport weights; the Gaussian covariance formulas of v1 do not hold for these transported laws.

Explicitly, model A is $z_t=x_t$; model B is $z_t=R_N(\sqrt\lambda\,T+\sqrt{1-\lambda}\,x_t)$; and model E is $z_t=R_N(s_ix_0+(1-s_i)x_t)$ with $s_i\sim\mathrm{Bernoulli}(\pi)$, so stayers retain their initial position. All innovations $\epsilon_t$ are standard normal. Model C shares A's latent law and alters only the response channel described above. In model D, the state-dependent coefficient $\rho(z)$ replaces the constant autoregressive coefficient inside the transported update. These equations and the numerical fitting rules specify the alternatives actually compared; a gap from the declared screen or solver does not establish nonexistence of a fit in a broader family.

% M58 (moved from Methods)
The 24-step scenario keeps the per-step coefficient 0.85, so without influence the per-wave latent autoregression falls to roughly $0.85^{24}\approx0.02$ before transport; it changes the annual dynamics as a whole, not only their temporal resolution. Under full anchoring ($\omega=1$) latent states never change regardless of $a$. A unit at $\omega=1$, whether complete and valid or complete but invalid, therefore carries no information about $a$ and is read only as a check of that degeneracy.

% M59 M08 M22 (moved from Methods)
The conditioning channel lets the signal share rise linearly with the number of earlier interviews, up to 0.95. This is a sensitivity of the response channel, not a model of panel conditioning, whose documented effect also includes true stabilization \cite{kraemer2025}. In the heteroscedastic channel, used only in the persistence arm of comparisons against C, the error multiplier $1+2e^{-z^2/2}$ is divided by its empirical root mean square over all $11\times N$ person-waves of the panel, non-respondents included, rather than by its expectation under a standard normal. Latent, error-free ordinal, latent-rank and mobility channels retain separate interpretations.

\subsection{Panel classifier and AUC interpretation}\label{sec:dd-classifier}

The fixed linear discriminant uses the pooled within-person correlations at lags three, four and five, with the ordinal categories treated as equally spaced scores. Standardization uses only training data; priors are equal and the covariance is $0.8S+0.2I$, where $S$ is the pooled within-class covariance of standardized training features. Training uses 2,000 panels per arm. Evaluation uses 4,000 new panels per arm, arranged as 2,000 networks with two panels each. No threshold or tuning search is performed.

For a given matched model and trained rule, under an equal-variance normal model for panel scores, $\mathrm{AUC}=\Phi(\delta/\sqrt2)$, where $\delta$ is the difference in mean scores in standard deviations. A one-sided test at level 0.05 that compares one panel with a known persistence score distribution then has power $\Phi(\delta-1.645)$. Article Table~\ref{main-tab:auc-translation} reports this translation for the three AUC targets. This illustrative translation is not an empirically validated model of the scores, the average power over the matching and training procedure, or the measured power of a test that fits persistence to one observed survey. AUC is the probability that an independently drawn influence panel ranks above an independently drawn persistence panel, with ties weighted by one half; it is neither classification accuracy at a fixed threshold nor the probability that influence is present in an observed panel.

\subsection{Feasibility screen and gap types}\label{sec:dd-feasibility}

% M10 M11 M12 M13 (moved from Methods)
Before classifier training, a declared feasibility mask compares the mean observed target lag-one correlation of the matching panels with the channel's continuous-scale upper bound. The ordinal rank map can change that bound, so the mask is a conservative operational screen, not an ordinal impossibility theorem. A cell that fails the screen receives status \texttt{family-limit-upper} (\texttt{family-limit-lower} for a negative target) and zero persistence-model evaluations; neither the solver nor validation runs. These cells mark targets beyond the continuous-scale bound at that signal share; they are gaps, not solver failures, and are reported separately from solver and validation gaps. The profile solver can also end in a family-limit status along $\eta$, but only after persistence-model evaluations. The evaluation count therefore distinguishes the two cases. An excluded, unmatchable or numerically failed cell is a gap, never counted as non-attainment. The full grid remains fixed.

% SUP02
For independent response error, continuous lag correlation is $\sqrt{r_t r_s}$ times latent lag correlation, giving the operational $r$ bound at constant signal share. Persistent style adds $\kappa\sqrt{(1-r_t)(1-r_s)}$ to the bound; conditioning uses the maximum pairwise bound. Cells on the bound remain eligible. Latent and error-free controls use bound one. These are screening rules before ordinal conversion, not empirical reliability estimates or proofs about the ordinal correlations.

\subsection{Conditional mean of the influence term under independence}\label{sec:dd-drift-mean}\label{supp:drift-mean}

Section~\ref{main-sec:methods-identification} of the article writes the influence term, to first order, as a conditional mean $m(z)=\mathbb E f(Z-z)$, with $Z\sim N(0,1)$, plus a deviation contributed by the particular contacts.
% M61 (moved from Methods)
With $g=Z-z$, the mean splits into attraction to contacts within distance one and repulsion from the others, $m(z)=\mathbb E[g\,\mathbf 1\{|g|\le1\}]-b\,\mathbb E[g\,\mathbf 1\{|g|>1\}]$, which at $b=0.25$ equals $1.25\,\mathbb E[g\,\mathbf 1\{|g|\le1\}]+0.25\,z$; for $z>0$ the attraction term pulls the person toward the population center and the repulsion term pushes the person away from it. Under this independent-normal-contact approximation, $m$ is odd and increases with $z$ at $b=0.25$, so its pre-standardization mean points away from the center for $z\neq0$: a push that is small near the center and grows in the tails. This is an analytic property of the approximation, not a sign or monotonicity guarantee for the finite network, the standardized influence term or the transported update. The attraction-only case has a mean directed toward the center under the same approximation.

% SUP15
Under the independence approximation in Section~\ref{main-sec:methods-identification} of the article, let $Z\sim N(0,1)$ and let $\varphi$ and $\Phi$ denote the standard-normal density and distribution function. The symbol $\varphi$ is distinct from the persistence parameter $\phi$. Integrating the frozen kernel at $b=0.25$ gives
\begin{align*}
m(z)&=1.25\bigl\{\varphi(z-1)-\varphi(z+1)\\
    &\quad-z[\Phi(z+1)-\Phi(z-1)]\bigr\}+0.25z,\\
m'(z)&=0.25+1.25\bigl\{\varphi(z-1)+\varphi(z+1)\\
     &\quad-[\Phi(z+1)-\Phi(z-1)]\bigr\}.
\end{align*}
Set $q(z)=\Phi(z+1)-\Phi(z-1)-\varphi(z-1)-\varphi(z+1)$, so $m'(z)=0.25-1.25q(z)$. The function is even, and
\begin{equation*}
q'(z)=2\varphi(z)e^{-1/2}\cosh z\,[z-2\tanh z].
\end{equation*}
For $z>0$, $h(z)=z-2\tanh z$ has $h(0)=0$, $h'(0)=-1$ and $h''(z)=4\,\mathrm{sech}^2z\,\tanh z>0$, with $h(z)\to\infty$. Thus $q$ first decreases and then increases to zero; its global maximum is $q(0)=2\Phi(1)-1-2\varphi(1)>0$. Consequently
\begin{align*}
m'(z)&\ge m'(0)=0.25-1.25\,c,\\
c&=2\Phi(1)-1-2\varphi(1),\\
m'(0)&\approx0.001565>0.
\end{align*}
The sign also follows from an explicit series bound. With $T_5=\sum_{k=0}^{5}(-1)^k/(2^k k!)$ and $S_4=\sum_{k=0}^{4}(-1)^k/[2^k k!(2k+1)]$, alternating-series bounds give $e^{-1/2}\ge T_5$ and $\int_0^1e^{-x^2/2}\,dx\le S_4$. Hence
\begin{equation*}
m'(0)\ge0.25+2.5\,(T_5-S_4)/\sqrt{2\pi}>0.0015.
\end{equation*}
Symmetry gives $m(-z)=-m(z)$ and $m(0)=0$, proving $m(z)>0$ for every $z>0$ and increasing magnitude with $|z|$. This describes the conditional mean before standardization and rank transport under independence, not an exact conditional law for the finite network.

Before standardization, the formula gives $m(0.5)\approx0.013$, $m(1)\approx0.085$ and $m(2)\approx0.40$. With attraction alone ($b=0$), it gives $m(1)\approx-0.13$. Under independence, the mean share of contacts within distance one is $P(|N(0,2)|\le1)\approx0.52$; transport restores the marginal but does not establish independence between a person and their finite-network contacts. % M61 SUP15

% Added in the journal revision (2026-10-13): analytic values of the frozen
% kernel, obtained by numerical integration; not a simulation result.
Under the same approximation, the deviation of the influence term from its conditional mean, contributed by the particular eight contacts, has variance $\mathrm{Var}\,f(Z-z)/8$ before standardization. Numerical integration of the frozen kernel at $b=0.25$ gives 0.0311 at $z=0$ and 0.0331 at $|z|=1$, with values in between for $|z|<1$. For $|z|\le1$, where $|m(z)|\le0.085$, its standard deviation is therefore about 0.18. Because a person's contacts are fixed and their states persist from step to step, this deviation is serially correlated within a person rather than independent across steps; the approximation describes its size, not its dependence. These values quantify the approximation discussed in Section~\ref{main-sec:methods-identification} of the article.

% Added in round 3 of the journal revision (2026-10-13): population integrals of
% the frozen kernel under the same approximation; analytic values, not results.
Read after the results became available, and added at the suggestion of an AI methods review of the journal version (Section~\ref{sec:ai}): over the whole population, with the person's own state also standard normal, numerical integration of the frozen kernel at $b=0.25$ gives $\mathrm{Var}\,m(z)\approx0.026$ for the conditional mean and $\mathbb E\,[\mathrm{Var}\,f(Z-z)]/8\approx0.031$ for the contact deviation, so the two parts are of similar size; about 98\% of $\mathrm{Var}\,m(z)$ comes from $|z|>1$. Most of the conditional mean is linear in the own state: its least-squares slope on $z$ is about 0.149, which accounts for about 84\% of $\mathrm{Var}\,m(z)$. A part linear in the own state acts like slightly stronger persistence, which matching to the lag-one and lag-two correlations can absorb (Section~\ref{main-sec:methods-identification} of the article explains why a linear dependence on the own state cannot be detected). Of the variance of the influence term before standardization, about 0.058, about 38\% is therefore linear in the own state, 7\% is the nonlinear, state-dependent part of the mean and 54\% is the contact deviation; with attraction alone ($b=0$) the shares are 23\%, 9\% and 67\%, and the linear part points toward the center. The exploratory variance-explained analysis, in which a spline of the own state absorbed little of the drift's contribution (Section~\ref{sec:x1-x2}), agrees with a small nonlinear own-state part. Like the values above, these describe the approximation, which ignores the correlation between a person and their contacts that builds up at positive strength.

If each person were instead pulled linearly toward the population mean, the influence term would be proportional to the person's own state, because transport keeps the population mean at zero; the update would remain an autoregression, and model A matched at lag one would have the same law, with population AUC 0.5 for any classifier. The protocol declared in advance that this class cannot be detected.

\subsection{Inherited design conventions}\label{sec:dd-conventions}

% SUP14
Table~\ref{tab:design-conventions} separates the role of three fixed constants from their recorded basis. The frozen v1 protocol declares them, and v4 retains them; no further rationale for their numerical values is recorded in those documents or the linked decisions.
\begin{table}[h]
\caption{Fixed conventions inherited from the original frozen protocol}\label{tab:design-conventions}
\begin{tabular}{@{}p{0.22\linewidth}p{0.30\linewidth}p{0.36\linewidth}@{}}
\toprule
Value & Role & Recorded basis and scope \\
\midrule
AUC 0.70 & Primary AUC target & Inherited operational convention, without universal substantive meaning; 0.60 and 0.80 remain secondary. \\
Shrinkage 0.2 & Fixed covariance $0.8S+0.2I$ & Inherited discriminator convention; the procedure includes no AUC-based hyperparameter search. \\
$\phi,\rho_0\in[0,0.999]$, $\eta\in[0,0.8]$ & Nonnegative parameter box of the persistence models & Inherited domain of the contrast, not an exhaustive characterization of the mathematical persistence families. \\
\botrule
\end{tabular}
\end{table}

\subsection{Matching: profile fit, validation and retry}\label{sec:dd-matching}

% M16 M17 M18 M19 M20 M21 (moved from Methods)
Each persistence model is matched to observed within-person lag correlations pooled over valid pairs and averaged across panels. B, D and E profile the second parameter $\eta\in\{\lambda,\gamma,\pi\}$: solve lag one in persistence $\phi$ at fixed $\eta$, then lag two along that profile. A, C and control refits solve lag one; lag two remains diagnostic. Matching and its independent validation pair target and persistence-model state, mask and error streams with common random numbers; training, testing and reference use independent classes and disjoint phases. Each fit attempt's profile search executes at most 96 persistence-model evaluations, including its Jacobian probes. Cached points are not recounted; validation and finite-difference stencils are counted separately. An accepted interior solution needs a real solver step and $|\bar g_k|\leq\tau_k=\min(5\times10^{-5},\max(10^{-6},0.25\,\mathrm{SE}_{\mathrm{fit},k}))$ for every fitted moment.

Before searching upward from the lower end of the lag-one-solvable profile ($\eta=0$, or the smallest $\eta$ at which lag one is solvable at $\phi=0.999$), the solver compares the lag-two residual there with its value 0.01 higher in $\eta$: a residual that does not rise, even after both inner roots are refined, gives \texttt{ill-conditioned}, and a rising residual with $\bar g_2\geq-\tau_2$ at the lower end gives \texttt{boundary}. Otherwise the solver accepts the first profile point it reaches at which every fitted moment meets tolerance. Monotonicity of the lag-two residual in $\eta$ is checked only at evaluated profile points, so on a hump-shaped profile the solver can accept a genuine interior root before it samples the region where the residual falls (search steps in Section~\ref{sec:odd}, Submodel~8).

Every accepted fit is checked on 2,048 fresh coupled pairs against $4\sqrt{\mathrm{SE}_{\mathrm{fit},k}^2+\mathrm{SE}_{\mathrm{val},k}^2}$ for every fitted moment. For a profiled B, D or E \texttt{boundary} fit (not a control refit) this includes lag two, so a constrained optimum whose lag-two mismatch exceeds four combined standard errors fails validation. One failed validation permits a prespecified 4,096-pair fit retry with its own 96-evaluation budget and fresh validation; a solver gap is not retried. A second failure is a gap.

% SUP03
The matching target is the mean across full-size panels of each pooled within-person lag correlation. The first attempt has 1,024 paired panels. Parameters satisfy $\phi,\rho_0\in[0,0.999]$ and $\eta\in[0,0.8]$. The profile prioritizes lag one. An anchor Jacobian seeds the outer bracket; guarded Illinois steps and bisection handle degenerate secants. Provisional inner roots are refined before acceptance and before declaring a non-monotone profile. Cached evaluations do not consume the 96-call budget twice. The lower end of the lag-one-solvable profile is accepted as a boundary only when its lag-one root meets tolerance and a local evaluation at $\eta=\min(0.8,\eta_{lower}+0.01)$ supports the constrained profile direction. An endpoint accepted without a real step is labeled \texttt{boundary}. Accepted statuses are \texttt{anchor-exact}, \texttt{interior} and \texttt{boundary}; \texttt{no-L1-root}, \texttt{non-monotone-profile}, \texttt{ill-conditioned}, \texttt{numerical-limit}, \texttt{family-limit-lower/upper} and \texttt{budget} are gaps. An exact anchor requires equality of all paired differences, not cancellation of their mean.

% SUP04
The independent validation uses 2,048 new pairs at the fitted parameter. Its rule includes both fitting and validation uncertainty; the 4,096-pair retry has a fresh 2,048-pair validation and its own 96-call budget. For the initial fit, $\mathrm{SE}_{val}/\mathrm{SE}_{fit}$ is about 0.71 under the declared pair counts; for the retry it is about 1.41. Each accepted fit also reports an equivalence statistic $|res|+2.2414\mathrm{SE}_{val}$ against 0.0003 and a local feature-amplification diagnostic. Neither diagnostic overrides the acceptance rule. Controls retain the second parameter from the corresponding observed fit in each retry and nuisance draw.

\subsection{Scope of matching precision and residual tolerance}\label{sec:dd-matching-scope}\label{supp:matching-scope}

% M66 (moved from Methods)
Matching estimates the influence arm's population moments from 1,024 (or 4,096) simulated panels, rather than targeting one realization. The solver's standard error pertains to coupled persistence-minus-target differences, not to the target mean alone. Their coupling covariance prevents reading the panel count as a fixed precision gain over fitting one observed panel.

For iid target panels, the standard deviation of one panel's lag statistic is $\sqrt{1024}=32$ times the standard error of its 1,024-panel mean. This identity concerns repeated target panels with the same variance, not the solver's residual standard error or the accuracy of fitting an empirical panel. The code forms paired differences $g=L_{\mathrm{persistence}}-L_{\mathrm{target}}$ and uses their sample standard deviation divided by the square root of the pair count. Their variance is $\mathrm{Var}(L_{\mathrm{persistence}})+\mathrm{Var}(L_{\mathrm{target}})-2\mathrm{Cov}(L_{\mathrm{persistence}},L_{\mathrm{target}})$, so common random numbers matter. % M66

% M69 (first sentence moved from Methods)
The target depends on the procedure: a cell that needs the retry is evaluated under a 4,096-panel matching distribution rather than 1,024, and the expectation is conditional on successful validated fits. The tolerance limits the fitting sample's mean residual in equations that satisfy the root criterion; it does not bound population mismatch or estimator bias. Unless the $10^{-6}$ floor binds, the limit is at most $0.25\,\mathrm{SE}_{\mathrm{fit}}$. A profiled B, D or E boundary may instead satisfy the lower-face direction rule with a larger lag-two residual and must then pass fresh validation at four combined standard errors. The centered nuisance variance is a variance of local AUC draws, not a mean squared error; it does not by itself cover systematic mismatch. Equivalence and amplification statistics remain diagnostics. % M69

\subsection{Local uncertainty, band construction and inversion}\label{sec:dd-band}

% M26 M27 (moved from Methods)
The score component uses 2,000 network-cluster bootstrap draws on fixed test scores, keeping two nested events together. Matching and training uncertainty uses 2,000 local draws: resample paired matching differences \emph{without centering}, solve and project a linearized parameter shift according to the profile's boundary rule, shift persistence-model training and reference features by a common-random-number finite-difference Jacobian of their mean, resample training networks by class, refit the LDA and evaluate a fixed reference. The nuisance variance is the centered variance of these AUCs, without division by draw count; exact-anchor fits have zero nuisance variance analytically. The two variance components add, $\mathrm{SE}^2=\mathrm{SE}_{\mathrm{scores}}^2+V_{\mathrm{nuis}}$. The prespecified finite-difference step is 0.01; the positive-strength dispersion check (G2) also checks 0.005.

% M28 M29 M30 (moved from Methods)
For each $r$, the band stacks the available strengths. Let $\epsilon=1/(2n)$ with $n=2{,}000$ networks per class. Each of the 2,000 replicates is centered on the point estimate, $p^*_b=\mathrm{clip}_{[0,1]}(\hat p+A_b-\bar A+\delta_b)$, where $A_b$ is the score-bootstrap AUC, $\bar A$ its mean over draws and $\delta_b$ the centered nuisance deviation of local draw $b$, with the same draw index at every strength. Estimates and replicates are then clipped to $[\epsilon,1-\epsilon]$ for the logit transform. The studentized deviation uses the logit standard error evaluated at each replicate, $t_b(a)=|\mathrm{logit}\,p^*_b-\mathrm{logit}\,\hat p|\,p^*_b(1-p^*_b)/\max(\mathrm{SE},\epsilon)$, not the plug-in value at $\hat p$. The critical value $c$ is the larger of 2.7729 and the empirical 0.95 quantile of $\max_a t_b(a)$ under the \texttt{higher} rule, which takes the order statistic at or above the 0.95 position without interpolation. The band is $\mathrm{logit}^{-1}\bigl(\mathrm{logit}\,\hat p\pm c\,\max(\mathrm{SE},\epsilon)/(\hat p(1-\hat p))\bigr)$, keeping the endpoint 0 or 1 when $\hat p$ equals it. These three readings (point-centered replicates, standard error evaluated at each replicate and the \texttt{higher} quantile) follow the frozen v1 code: the first is in its score-uncertainty function and the other two in its band function. The v4 band re-implements them with the total standard error and fixes the floor at 2.7729 rather than recomputing it from the number of available strengths. A plug-in reading of the same specification gives critical values up to 2.1 times larger (Section~\ref{sec:verification-docking}), so the band width depends on this implementation choice.

% M31 M73 M74 (moved from Methods)
Each band is simultaneous over available strengths within one comparison/$r$ unit. No joint coverage is declared across values of $r$, persistence models or channels. Because seeds also differ across values of $r$ in the confirmatory run (Section~\ref{sec:dd-seeds}), statements about the shape of the curve across $r$, such as monotonicity or the range of $r$ over which a boundary is attained, are descriptive; no 81-cell band is prespecified. A strict one-sided crossing claim needs a lower band above 0.70; inclusive band inversion bounds the discrete first and sustained boundaries. The reported boundaries and intervals follow this inclusive inversion, and the word ``crossing'' in a claim is reserved for strengths whose lower band exceeds 0.70.

% M72 M24 M25 (moved from Methods)
The grid is coarse above 0.10, and intervals spanning several grid points are expected: in the v1 precision planning, which is not a power statement, the probability of an interval that covers the true boundary and spans at most one grid step was 43.9\% or 19.2\%, for nuisance standard deviations of 0.01 or 0.02, when the AUCs bracketing the boundary are 0.66 and 0.74. Neither interpolation nor monotone smoothing is allowed; gaps, right censoring and boundaries at the first grid point remain visible. Neither the solver nor the band enforces monotonicity in $a$. The sustained boundary accounts for any reversal of the estimated AUC across strengths, whatever its source, by requiring the estimate at every higher available strength to reach the target. Boundaries are computed on the available strengths of a curve. The first boundary and each end of its interval are retained only when strictly below the first gap; the sustained boundary and each end of its interval only when strictly above the last gap; otherwise the entry is censored and labeled with that gap's strength. If the sustained boundary on the available strengths lies below the last gap, it is censored rather than moved to the first strength above that gap; a sustained boundary at the first available strength above the last gap is retained, although the unobserved gap strength could also have met the target.

% SUP05
For each local draw, resampled matching differences are averaged without centering, then the fitted Jacobian determines a parameter perturbation. Projection onto the box re-solves lag one after projecting $\eta$ and reproduces the boundary decision. The mean-feature Jacobian $D$ shifts persistence-model training and reference lag features by $D\Delta\theta$. Common-random-number stencils on matching panels use a step of 0.01, inward at an upper face. G2 alone repeats at 0.005 and evaluates a covariance-Jacobian approximation as a diagnostic; an approximate AUC correction above 10\% of $\sqrt{V_{nuis}}$ is flagged, not silently incorporated. Training networks are resampled separately by class with both events together. The fixed reference has 2,000 networks and two events per class. The local variance is centered; mean shift, fourth-moment Monte Carlo standard error and boundary-projection frequency are diagnostics. Anchor-exact draws set $\Delta\theta=0$ and $V_{nuis}=0$ analytically.

% SUP06
The score bootstrap uses 2,000 draws and aligned network indices across strengths. Local draws likewise share indices across strengths in one curve at fixed $r$. The logit transform clips probabilities only for transformation to $[1/(2n),1-1/(2n)]$; an estimated standard error has the same floor before logit scaling. Raw zero and one band endpoints remain reportable. Inversion uses the first upper-band crossing for the lower grid bound and the first lower-band crossing for the upper bound, including equality. A gap prevents declaring a minimum across it. The $a=0$ band must contain 0.5; its failure invalidates the comparison at that value of $r$. % SUP06

\paragraph{Scope of these approximations.}
Added in the journal revision, after the results became available, at the suggestion of an AI statistical review (Section~\ref{sec:ai}). The sum $\mathrm{SE}^2=\mathrm{SE}_{\mathrm{scores}}^2+V_{\mathrm{nuis}}$ is a first-order approximation that treats evaluation as independent of fitting and training, not an exact decomposition: the score variance is estimated at one fitted rule, and the local propagation shifts feature means instead of regenerating the whole fitted procedure with its retries, validation selection and fit failures. The floor 2.7729 is the two-sided normal Bonferroni critical value for nine strengths; it protects against multiplicity only if the marginal logit pivots are adequately calibrated, and it corrects neither bias nor an inaccurate standard error. The bands are therefore approximate simultaneous bands with a normal Bonferroni floor, not distribution-free guarantees. Given simultaneous coverage on a complete grid, inverting a band bounds the first and sustained boundaries even where the AUC is not monotone in $a$; a sustained boundary retained at the first strength above the last gap is only an upper bound on the sustained boundary of the full grid.

For example, at $r=\resultnum{v4.tab.T1.B.r055.r}{0.55}$ the fits against B failed at $a=\resultnum{j.b.r055.sustained.gap.0.a}{0.60}$ and $\resultnum{j.b.r055.sustained.gap.1.a}{0.90}$ (article, Table~\ref{main-tab:boundaries-B}): the first boundary, $a=\resultnum{j.b.r055.first.point}{0.10}$, lies below the first gap and is reported, but no strength lies above the last gap, so no sustained boundary is declared; the rule does not assume that the AUC stayed above 0.70 where no fit exists.

\subsection{Seeds}\label{sec:dd-seeds}

% M32 M33 M34 M35 (moved from Methods)
Seed material is the string \texttt{namespace|block|phase|role|index}; the first eight bytes of its SHA-256 digest, read big-endian, seed a NumPy generator. The zero-strength fitting check (G0) uses one calibration seed plan for all 81 cells, so random numbers are common across $a$ and $r$, as in v1. In the confirmatory run, each value of $r$ receives its own confirmatory seed block. All strengths and all comparison units at one $r$ therefore share state, mask, error, training, test, reference, nuisance and bootstrap streams, while different values of $r$ use separate streams. A stencil namespace is declared and included in the seed-collision audit, but no panel is drawn from it. In the confirmatory run, finite-difference stencils re-evaluate the matching-phase streams at shifted parameters, and the covariance diagnostic of the positive-strength dispersion check (G2) re-evaluates the training and reference streams in the same way, preserving common random numbers. The calibration, gate, nuisance, bootstrap, validation and stencil namespaces are separate from the reserved confirmatory namespaces. The seed audit includes all gate plans and derived namespaces, and phase manifests must bind code, environment, host and seed material. % SUP11

\subsection{Inventory tiers and the reading across persistence models}\label{sec:dd-inventory}

% M42 M43 M79 (moved from Methods)
The machine-readable inventory retains all declared comparisons. Fixed tiers are: the main curve, B on the main grid (tier 1); the secondary persistence models A, C, D and E (tier 2); three measurement controls and mobility against applicable persistence models (tier 3); response-style and conditioning scenarios (tier 4); attraction-only and anchoring scenarios (tier 5); and the 24-step scenario (tier 6). Tier 1 supports a claim only about B. Cross-mechanism interpretations require complete valid grids for the relevant secondary persistence models. When tier 2 is complete and valid and the main B curve is valid, the boundaries against B and D are reported side by side. Under these conditions, if B crosses while D's valid upper band stays below the target, the result is compatible with non-Gaussian persistence but gives no specific evidence of influence. If D's band spans the target, the reading is inconclusive. If tier 2 is not entirely complete and valid, no comparison across mechanisms is interpreted. This compares states obtained by separate band rules; it is neither a test of the difference between B and D AUCs nor a comparison with declared joint 95\% confidence.

The execution order, the release rule and the operational amendment of 1 October 2026 (v4.4) are in Section~\ref{sec:record-execution}; the unit statuses and the inventory are in Section~\ref{sec:inventory}.

For models A, B, D and E the zero-strength check concerns an identity, so it checks, in one draw per unit, the score part of the band at zero strength, not the matching at positive strength.
